%% kapa: exam=1 round=36 year=2025 period=2 form=A subject=appraisal_law subject_ko=감정평가관계법규 seq=32 number=32
%% source: Q-Net 감정평가사 기출문제 (cst00309) artlSeq=5231525 file=raw/1008/5231525/2025년 제36회 감정평가사 1차시험 2교시 문제지 원본.pdf
%% conversion: kordoc 4.24.51 -> tools/gen1.py (auto); review=proofread-sonnet; flags=none
\begin{question}{32}
부동산등기법령상 등기의 관할에 관한 설명으로 옳지 않은 것은?
\begin{choices}
  \item 지방법원장은 어느 등기소의 관할에 속하는 사무를 다른 등기소에 위임하게 할 수 있다.
  \item 관할 등기소가 다른 여러 개의 부동산과 관련하여 등기목적과 등기원인이 동일한 경우 그 중 하나의 관할 등기소에서 등기사무를 담당할 수 있다.
  \item 상속으로 인한 등기신청의 경우 부동산의 관할 등기소가 아닌 등기소도 그 신청에 따른 등기사무를 담당할 수 있다.
  \item 사건이 그 등기소의 관할이 아닌 경우의 등기신청은 각하사유에 해당한다.
  \item 부동산이 여러 등기소의 관할구역에 걸쳐 있을 때에는 대법원규칙으로 정하는 바에 따라 각 등기소를 관할하는 상급법원의 장이 관할 등기소를 지정한다.
\end{choices}
\end{question}
\answer{1}
